%-----------------------------------------------------------------------------------------------
\makeatletter
\immediate\write18{datelog > \jobname.info} % site script for $(date '+%Y-%m-%d %Hh%Mm%Ss')
\makeatother
%-----------------------------------------------------------------------------------------------
%-----------------------------------------------------------------------------------------------
\usetheme{Copenhagen}
\usepackage{beamercolorthemeUTF2}
\usefonttheme{serif}
%-----------------------------------------------------------------------------------------------
\usepackage[utf8]{inputenc}
\usepackage[greek,french,english,brazil]{babel} % last becomes the active one
\usepackage{pslatex}
\usepackage{amssymb,amsmath}
\usepackage{soul}
\usepackage[squaren,Gray,cdot]{SIunits}
\usepackage[nice]{nicefrac}
\usepackage{tikz}
\usepackage{amscd}
\usepackage{stmaryrd}
\usepackage{scalerel}
\usepackage{xspace}
%-----------------------------------------------------------------------------------------------


%-----------------------------------------------------------------------------------------------
%-----------------------------------------------------------------------------------------------
% Mathematical
%-----------------------------------------------------------------------------------------------
\newcommand{\vet}[1]{\underline{{#1}}}
\newcommand{\mat}[1]{\underline{\underline{{#1}}}}
\newcommand{\cub}[1]{\underline{\underline{\underline{{#1}}}}}
\newcommand{\eqdef}{{\ensuremath\stackrel{\text{\tiny def}}{=}}}
%-----------------------------------------------------------------------------------------------
% Linguistic
%-----------------------------------------------------------------------------------------------
\newcommand{\GRtxt}[1]{\begin{otherlanguage}{greek}{{#1}}\end{otherlanguage}}
\newcommand{\FRtxt}[1]{\begin{otherlanguage}{french}{{#1}}\end{otherlanguage}}
%-----------------------------------------------------------------------------------------------
% Presentation
%-----------------------------------------------------------------------------------------------
\newcommand{\BkgImgH}[1]{% Places an image centered on the slide background filling the height
    \usebackgroundtemplate{\parbox{\paperwidth}{%
        \vspace*{1sp}\centering\includegraphics[height=\paperheight]{{#1}}
}}}
\newcommand{\BkgImgW}[1]{% Places an image centered on the slide background filling the width
    \usebackgroundtemplate{\parbox{\paperwidth}{%
        \vspace*{1sp}\centering\includegraphics[width=\paperwidth]{{#1}}
}}}
\newcommand{\ArtEndH}[3]{% Transitions to plain image (last) slide: #1:prefix #2,#3:extensions
    \BkgImgH{root/../art/#1.#2}
    \frame<handout:0>[plain]{%
        \transdissolve\vspace*{72mm}\color{white}\scriptsize\bf\input{root/../art/#1.#3}}
    \usebackgroundtemplate{\mbox{~}}
}
\newcommand{\ArtEndW}[3]{% Transitions to plain image (last) slide: #1:prefix #2,#3:extensions
    \BkgImgW{root/../art/#1.#2}
    \frame<handout:0>[plain]{%
        \transdissolve\vspace*{72mm}\color{white}\scriptsize\bf\input{root/../art/#1.#3}}
    \usebackgroundtemplate{\mbox{~}}
}
\newcommand{\ImgColW}[3]{% Inserts a full-width image in a column
    \includegraphics[width=\columnwidth]{root/../art/#1.#2}\\[-0.5\baselineskip]
    \parbox{\columnwidth}{\tiny\hfill\scalebox{0.85}{\input{root/../art/#1.#3}}}
}
\newcommand{\txtpic}[1]{%
    \fcolorbox{lightgray}{white!90!black}{{#1}} 
}
%-----------------------------------------------------------------------------------------------


%-----------------------------------------------------------------------------------------------
\title{A.03.03 -- Balanço de Energia}
\subtitle{(Sistemas Fechados)}
\author{Prof.~C.~Naaktgeboren, PhD}
\date{{\scriptsize\tt%
    \includegraphics[height=6.0mm]{root/00-res/cc/by-nc-nd-88x31.pdf}\\[\smallskipamount]
    https://github.com/CNThermSci/ApplThermSci\\
    Compiled on \input{\jobname.info}
}}
%-----------------------------------------------------------------------------------------------
\begin{document}
%-----------------------------------------------------------------------------------------------
\logo{%
    \parbox{158mm}{% There's a 1mm gap on each side of the 160mm x 90mm slide logo line
        \includegraphics[height=6.5mm]{root/00-res/UTFPR/UTFPR-logo-A.pdf}\hfill%
        \includegraphics[height=6.5mm]{root/00-res/logo/CNThermSci-logo-A.pdf}%
%   \parbox{126mm}{% There's a 1mm gap on each side of the 128mm x 96mm slide logo line
%       \includegraphics[height=6.0mm]{root/00-res/UTFPR/UTFPR-logo-A.pdf}\hfill%
%       \includegraphics[height=6.0mm]{root/00-res/logo/CNThermSci-logo-A.pdf}%
}} % Alpha logos
%-----------------------------------------------------------------------------------------------
\frame{\titlepage}
%-----------------------------------------------------------------------------------------------

%-----------------------------------------------------------------------------------------------
\frame{\tableofcontents}
%-----------------------------------------------------------------------------------------------

%-----------------------------------------------------------------------------------------------
\section{Balanço de Energia}
%-----------------------------------------------------------------------------------------------

%-----------------------------------------------------------------------------------------------
\subsection{Primeira Lei da Termodinâmica}
%-----------------------------------------------------------------------------------------------

    % !j 96 -i8
    %-------------------------------------------------------------------------------------------
    \begin{frame}{Enunciado}\vspace*{-2em}
        \uncover<1->{A \alert{1$^{\mathrm{a}}$ lei da Termodinâmica} estabelece que:}
        \begin{itemize}
            \item<1-> \alert{Energia} é uma quantidade \alert{conservada}. \\[\bigskipamount]
        \end{itemize}
        \uncover<2->{Este \alert{princípio da conservação da energia}:}
        \begin{itemize}
            \item<2-> É exaustivamente \alert{confirmado} em experimentos.%\\[\bigskipamount]
        \end{itemize}
    \end{frame}
    %-------------------------------------------------------------------------------------------

    % !j 96 -i8
    %-------------------------------------------------------------------------------------------
    \begin{frame}{Algumas Implicações}\vspace*{-2em}
        \uncover<1->{Logo, no universo observável:}
        \begin{itemize}
            \item<1-> \alert{Não há processos físicos} que \alert{criem} energia,
                \\[\medskipamount]
            \item<2-> \alert{Nem processos físicos} que \alert{destruam} energia.
                \\[\medskipamount]
            \item<3-> \alert{Processos físicos} podem apenas \alert{converter} energia de uma
                forma a outra. \\[\medskipamount]
        \end{itemize}
        \uncover<4->{A Relatividade Especial de Einstein:}
        \begin{itemize}
            \item<4-> \alert{Unificou} as conservações de \alert{massa} e de \alert{energia};
                \\[\medskipamount]
            \item<5-> Através da \alert{equivalência massa-energia} expressa por $E_{eq} =
                c^2m$.  \\[\medskipamount]
            \item<6-> Assim, a quantidade $E_{tot} = c^2m + E_{outras}$ do universo é
                \alert{conservada}. %\\[\medskipamount]
        \end{itemize}
    \end{frame}
    %-------------------------------------------------------------------------------------------

    % !j 96 -i8
    %-------------------------------------------------------------------------------------------
    \begin{frame}{Aplicações}\vspace*{-2em}
        \begin{columns}
        \column{0.50\textwidth}
        \uncover<1->{A 1$^{\mathrm{a}}$ lei é \alert{central} em Termodinâmica.\par\noindent
            Suas aplicações são \alert{vastas} e incluem:}
        \begin{itemize}
            \item<1-> Princípio em variedade de \alert{deduções}; \\[\medskipamount]
            \item<2-> Instrumental na \alert{definição} de \alert{propriedades}.
                \\[\medskipamount]
            \item<3-> Cálculos de \alert{processos} energéticos. %\\[\medskipamount]
        \end{itemize}
        \column{0.50\textwidth}
        \uncover<4->{Exemplo: O que é, afinal, ``energia''?}
        \begin{itemize}
            \item<6-> ``Energia é uma \alert{quantidade} (escalar) \\[\medskipamount]
            \item<7-> \hspace{1ex}que é \alert{conservada} na natureza \\[\medskipamount]
            \item<8-> \hspace{1ex}e que possui \alert{unidades de
                \kilogram\usk\squaren\meter\per\squaren\second}.'' %\\[\medskipamount]
        \end{itemize}
        \uncover<5->{\hfill --- Jack P.~Holman (SMU)}
        % https://www.smu.edu/News/2013/jack-holman-06may2013
        \end{columns}
    \end{frame}
    %-------------------------------------------------------------------------------------------

%-----------------------------------------------------------------------------------------------
\subsection{Balanço de Energia}
%-----------------------------------------------------------------------------------------------

    % !j 96 -i8
    %-------------------------------------------------------------------------------------------
    \begin{frame}{Balanço de Energia}\vspace*{-2em}
        \uncover<1->{A 1$^{\mathrm{a}}$ lei é matematicamente expressa por meio de
        \alert{balanço de energia}.}\\[\medskipamount]
        \uncover<2->{Em um \alert{processo}, o balanço de energia é dado por:}%
        \begin{equation*}%
            \uncover<3->{\left(\begin{array}{c}
                \text{Total de energia que} \\
                \text{\alert{entra} no sistema}
            \end{array}\right)-}
            \uncover<4->{\left(\begin{array}{c}
                \text{Total de energia} \\
                \text{que \alert{sai} do sistema}
            \end{array}\right)=}
            \uncover<5->{\left(\begin{array}{c}
                \text{Variação \alert{líquida} de} \\
                \text{energia no sistema}
            \end{array}\right),}%
        \end{equation*}%
        \uncover<6->{que matematicamente se escreve:}%
        \begin{equation*}%
            \uncover<6->{\alert{E_{ent} - E_{sai} = \Delta{E}_{sist} = E_2 -
            E_1},\quad\mbox{para um processo $1$--$2$.}}%
        \end{equation*}%
        \uncover<7->{Assim, se $E_1$, $E_{ent}$ e $E_{sai}$ são conhecidos, então: $E_2 = E_1 +
        E_{ent} - E_{sai}$.}
    \end{frame}
    %-------------------------------------------------------------------------------------------

%    % !j 96 -i8
%    %-------------------------------------------------------------------------------------------
%    \begin{frame}{Balanço de Energia -- Formas}\vspace*{-2em}
%        \begin{align*}
%            % --- first line
%            \uncover<1->{
%                \alt<1,2,4,8>%
%                {\alert{E_{ent} - E_{sai}} & \alert{= \Delta{E}_{sist} = E_2 - E_1}}%
%                {E_{ent} - E_{sai} & = \Delta{E}_{sist} = E_2 - E_1}%
%            } &
%            \uncover<2->{\quad\stackrel{(d/dt)}{\longrightarrow}\quad} & &
%            \uncover<3->{
%                \alt<3,6,8>%
%                {\alert{\dot{E}_{ent} - \dot{E}_{sai}} & \alert{= \left.\frac{dE}{dt}\right|_{sist}}}%
%                {\dot{E}_{ent} - \dot{E}_{sai} & = \left.\frac{dE}{dt}\right|_{sist}}%
%            } \\[\bigskipamount]
%            % --- second line
%            \uncover<4->{\downharpoonright & (\div m)} & & &
%            \uncover<6->{\downharpoonright & (\div m)} \\[\bigskipamount]
%            % --- third line
%            \uncover<5->{
%                \alt<5,6,8>%
%                {\alert{e_{ent} - e_{sai}} & \alert{= \Delta{e}_{sist} = e_2 - e_1}}%
%                {e_{ent} - e_{sai} & = \Delta{e}_{sist} = e_2 - e_1}%
%            } &
%            \uncover<6->{\quad\stackrel{(d/dt)}{\longrightarrow}\quad} & &
%            \uncover<7->{
%                \alt<7,8>%
%                {\alert{\dot{e}_{ent} - \dot{e}_{sai}} & \alert{= \left.\frac{de}{dt}\right|_{sist}}}%
%                {\dot{e}_{ent} - \dot{e}_{sai} & = \left.\frac{de}{dt}\right|_{sist}}%
%            } \\%[\medskipamount]
%        \end{align*}
%    \end{frame}
%    %-------------------------------------------------------------------------------------------

    % !j 96 -i8
    %-------------------------------------------------------------------------------------------
    \begin{frame}{Balanço de Energia -- Formas}\vspace*{-2em}
        \small
        \begin{align*}
            % --- `zeroth' line
            \uncover<1->{ & \text{Processo}} &
            \uncover<2->{
                \stackrel{\scriptstyle d()}{\longrightarrow}
            } & &
            \uncover<3->{\text{Diferenc}&\text{ial}} &
            \uncover<4->{
                \stackrel{\scriptstyle /dt}{\longrightarrow}
            } & &
            \uncover<5->{\text{Taxa} & }
            \\[\bigskipamount]
            % --- first line
            \uncover<1->{
                \alt<1,2,6,12>%
                {\alert{E_{ent} - E_{sai}} & \alert{= \Delta{E}_{sist} = E_2 - E_1}}%
                {E_{ent} - E_{sai} & = \Delta{E}_{sist} = E_2 - E_1}%
            } &
            \uncover<2->{
                \stackrel{\scriptstyle d()}{\longrightarrow}
            } & &
            \uncover<3->{
                \alt<3,4,8,12>%
                {\alert{\delta{E}_{ent} - \delta{E}_{sai}} & \alert{= dE_{sist}}}%
                {\delta{E}_{ent} - \delta{E}_{sai} & = dE_{sist}}%
            } &
            \uncover<4->{
                \stackrel{\scriptstyle /dt}{\longrightarrow}
            } & &
            \uncover<5->{
                \alt<5,10,12>%
                {\alert{\dot{E}_{ent} - \dot{E}_{sai}} & \alert{= \left.\frac{dE}{dt}\right|_{sist}}}%
                {\dot{E}_{ent} - \dot{E}_{sai} & = \left.\frac{dE}{dt}\right|_{sist}}%
            }
            \\[\bigskipamount]
            % --- second line
            \uncover<6->{\text{(int.)}\qquad & \quad\downharpoonright{\scriptstyle \div m}} & & &
            \uncover<8->{\downharpoonright {\scriptstyle \div m} & } & & &
            \uncover<10->{\downharpoonright {\scriptstyle \div m} & }
            \\[\bigskipamount]
            % --- third line
            \uncover<7->{
                \alt<7,8,12>%
                {\alert{e_{ent} - e_{sai}} & \alert{= \Delta{e}_{sist} = e_2 - e_1}}%
                {e_{ent} - e_{sai} & = \Delta{e}_{sist} = e_2 - e_1}%
            } &
            \uncover<8->{
                \stackrel{\scriptstyle d()}{\longrightarrow}
            } & &
            \uncover<9->{
                \alt<9,10,12>%
                {\alert{\delta{e}_{ent} - \delta{e}_{sai}} & \alert{= de_{sist}}}%
                {\delta{e}_{ent} - \delta{e}_{sai} & = de_{sist}}%
            } &
            \uncover<10->{
                \stackrel{\scriptstyle /dt}{\longrightarrow}
            } & &
            \uncover<11->{
                \alt<11,12>%
                {\alert{\dot{e}_{ent} - \dot{e}_{sai}} & \alert{= \left.\frac{de}{dt}\right|_{sist}}}%
                {\dot{e}_{ent} - \dot{e}_{sai} & = \left.\frac{de}{dt}\right|_{sist}}%
            }
        \end{align*}
    \end{frame}
    %-------------------------------------------------------------------------------------------

    % !j 96 -i8
    %-------------------------------------------------------------------------------------------
    \begin{frame}{Balanço de Energia -- $E_{ent}$, $E_{sai}$}\vspace*{-2em}
        \uncover<1->{Em \alert{sistemas compressíveis simples}, $E_{ent}$ e $E_{sai}$ podem ser
        apenas nas formas de:}\\[\medskipamount]
        \begin{enumerate}
            \item<2-> \alert{calor} e \\[\medskipamount]
            \item<3-> \alert{trabalho}. \\[\bigskipamount]
        \end{enumerate}
        \uncover<4->{Assim, no balanço de energia:}
        \begin{align*}%
            \uncover<5->{
                \alert{E_{ent}} & \alert{= Q_{ent} + W_{ent}},\qquad\text{e}
            } \\[\medskipamount]
            \uncover<6->{
                \alert{E_{sai}} & \alert{= Q_{sai} + W_{sai}}.
            } %\\[\medskipamount]
        \end{align*}%
    \end{frame}
    %-------------------------------------------------------------------------------------------

    % !j 96 -i8
    %-------------------------------------------------------------------------------------------
    \begin{frame}{Balanço de Energia -- $E_{sist}$}\vspace*{-2em}
        \uncover<1->{Em \alert{sistemas clássicos} (pré-relativísticos) \alert{não reativos},
        $E_{sist} = E_{micro} + E_{macro}$:}\\[\medskipamount]
        \begin{enumerate}
            \item<2-> $E_{micro} \equiv U_{sist}$, a \alert{energia interna}, em
                \kilo\joule\per\kilogram; \\[\medskipamount]
            \item<3-> $E_{macro} = E_c + E_p$, a soma das energias \alert{cinética}, e
                \alert{potencial}, em \kilo\joule\per\kilogram, onde: \\[\medskipamount]
            \item<4-> $E_c = m\mathbb{V}^2/2$, com $\llbracket\mathbb{V}\rrbracket =
                \sqrt{\kilo\joule\per\kilogram} = \unit{\sqrt{1000}}{\meter\per\second} \approx
                \unit{31.62}{\meter\per\second}$, ou \\[\medskipamount]
            \item<5-> $E_c = m{\scriptstyle\mathbb{V}}^2/2$, com
                $\llbracket{\scriptstyle\mathbb{V}}\rrbracket = \meter\per\second$;
                \\[\medskipamount]
            \item<6-> $E_p = g\mathbb{Z}$, com $\llbracket g\rrbracket =
                \meter\per\squaren\second$ e $\llbracket\mathbb{Z}\rrbracket = \kilo\meter$, ou
                \\[\medskipamount]
            \item<7-> $E_p = g{\scriptstyle\mathbb{Z}}/1000$, com $\llbracket g\rrbracket
                = \meter\per\squaren\second$ e
                $\llbracket{\scriptstyle\mathbb{Z}}\rrbracket = \meter$.
                %\\[\medskipamount]
        \end{enumerate}
    \end{frame}
    %-------------------------------------------------------------------------------------------

%-----------------------------------------------------------------------------------------------
\section{Tópicos de Leitura}
%-----------------------------------------------------------------------------------------------

    %------------------------------------------------------------------------------------------
    \begin{frame}[allowframebreaks]{Tópicos de Leitura}
        \begin{thebibliography}{Çengel, Y.~A., 2013}
            \bibitem[Çengel, Y.~A., 2013]{2013-CengelYA+BolesMA-AMGH}
                Çengel, Y.~A. e Boles, M.~A.
                \newblock{%
                    {\em Termodinâmica $7^\mathrm{a}\!$ Edição\/}.
                    \alert{Seções~2-6 e 4-2.}
                }
                \newblock{\footnotesize AMGH. Porto Alegre. ISBN 978-85-8055-200-3.}
        \end{thebibliography}
    \end{frame}
    %------------------------------------------------------------------------------------------

    % Finishes with stunning image, with credit
    \ArtEndW{italy-1587287_1280}{jpg}{txt}

%-----------------------------------------------------------------------------------------------
\end{document}
%-----------------------------------------------------------------------------------------------

